\subsection{A signed Galois--fermion cut for \texorpdfstring{$\RR{8}{1660}$}{R6\_8\_1660}}
\label{subsec:r1660-signed-galois}
The following necessary condition removes one further labelled profile from
 the frozen $1670$-profile frontier.  It leaves $1669$ profiles, in branches
of sizes $163,753,753$; it does not resolve categorifiability.

Work in the established pseudo-unitary Galois form, with
$\mathcal M=Z(\mathcal C)$, $T=\sigma X$, and restriction vectors in the
order $(1,g,X,gX,\sigma,T)$.  Let $h$ be the central fermion and let
$\widehat\alpha=Z(\alpha)$ be the previously constructed involution.
For $a\in(\mathbb Z/48\mathbb Z)^\times$, write $\pi_a$ for the Galois
permutation and $\delta_a(Z)\in\{\pm1\}$ for its sign, so that for the
normalized modular matrix $S$,
\[
 \sigma_a(S_{L,Z})=\delta_a(Z)S_{L,\pi_a Z},
 \qquad \theta_{\pi_a Z}=\theta_Z^{a^2}.
\]

\begin{lemma}[The signed restriction action]
For every central simple $Z$ with restriction vector $v$,
\[
 v_{\pi_5 Z}=\delta_5(Z)B_5v,\qquad
 B_5=\begin{pmatrix}
1&0&0&0&0&0\\
0&1&0&0&0&0\\
2&2&-1&-2&-2&2\\
2&2&-2&-1&-2&2\\
0&0&0&0&-1&0\\
2&2&-2&-2&-4&3
\end{pmatrix}.
\]
The sign is the one making the vector nonnegative.  Moreover $\pi_5$
commutes with tensoring by $h$ and with $\widehat\alpha$, and preserves
monodromy charge with $h$.
\end{lemma}
\begin{proof}
The six constituents $L_\chi$ of $I(1)$ correspond to the six characters
of $K_0(\mathcal C)$, with formal codegrees $f_\chi$.  The known induction
rows are
\[
 S_{L_\chi,Z}=\frac{\chi(FZ)}{f_\chi}.
\]
On $K=\mathbb Q(\sqrt2,\sqrt3)$, $\sigma_5$ negates both square roots.
For all six characters and all six basis elements, direct arithmetic gives
\[
 \sigma_5\!\left(\frac{\chi(v)}{f_\chi}\right)
       =\frac{\chi(B_5v)}{f_\chi}.
\]
The character table is invertible, proving the asserted restriction formula.
The verifier checks these $36$ identities exactly in the basis
$1,\sqrt2,\sqrt3,\sqrt6$, and checks $B_5^2=1$.

Write $\epsilon_Z$ for the charge.  Since
$S_{h,Z}/S_{1,Z}=\epsilon_Z\in\{\pm1\}$, the column Galois identity
shows that $\pi_5$ preserves charge.  Also
$S_{hY,Z}=\epsilon_ZS_{Y,Z}$ and $d_{hY}=d_Y$ imply
$\delta_5(hY)=\delta_5(Y)$; applying Galois to the row identity and
using invertibility of $S$ gives $\pi_5(hY)=h\pi_5(Y)$.
Finally, the permutation matrix of a braided autoequivalence commutes with
$S$ and is rational, so it commutes with the signed Galois matrix
$S^{-1}\sigma_5(S)$.  Preservation of dimensions also preserves its signs,
proving commutation with $\widehat\alpha$.
\end{proof}

\begin{lemma}[Four-divisibility of the spin fixed locus]
In a restriction fiber
\[
 v=(0,0,b,c,1,3),\qquad b+c=4,
\]
the number of $h$-fixed simples is divisible by four.
\end{lemma}
\begin{proof}
The signed formula gives
$B_5v=-(0,0,c,b,1,3)$, so
$\gamma=\pi_5\widehat\alpha$ preserves the fiber.
Its simples have primitive sixteenth-root twists; therefore
$\theta_{\gamma Z}=\theta_Z^{25}=-\theta_Z$.
The preceding lemma implies that $\gamma$ preserves the $h$-fixed locus.
Duality preserves that locus, the restriction vector and the twist.
It acts freely: a self-dual simple with this restriction would have to
have indicator $-1$ from the odd $\sigma$ multiplicity and indicator
$+1$ from the odd $T$ multiplicity.  Thus each twist occurs an even
number of times in the fixed locus.  The bijection $\gamma$ gives equal
multiplicities for the two twists $\theta,-\theta$, proving divisibility
by four.  No claim that $\gamma$ is a tensor functor is used.
\end{proof}

Here is the complete exclusion of the zero-based raw profile $133$ in
low frame zero, character labelling $(0,5,4,1,3,2)$.
Remove the four forced even low/high pairs.  The remaining $18$ high
simples have the following restrictions:
\[
\begin{array}{c|c|c}
\text{symbol}&\text{restriction}&\text{multiplicity}\\\hline
A&T&2\\
A'&2X+2gX+3T&2\\
B&X+gX+2T&2\\
C&4X+4gX+4T&2\\
\Sigma&2X+2gX+\sigma+3T&8\\
P&2X+2gX+2T&1\\
Q&4X+4gX+8T&1
\end{array}
\]
All these fibers are $h$-invariant.  Let $e_U,o_U,t_U$ be their numbers
of even $h$-pairs, odd $h$-pairs and odd $h$-fixed simples.  The signed
Galois map exchanges $A,A'$, so all three orbit counts agree for this
pair.  Put $a=e_A=e_{A'}$, $b=e_B$, $c=e_C$, and $z=e_\Sigma$.
Then $a,b,c\in\{0,1\}$, $e_P=e_Q=0$, and the previously proved duality
parity gives $z\in\{0,2,4\}$.  The even-dimension equation is exactly
\[
 3a+b+8c+3z=15.
\]
Its only solutions are
\[
 (a,b,c,z)=(1,0,0,4),\qquad(0,1,1,2).
\]
For total remaining orbit counts $E,O,H$, the already established rank
identity in this low frame is $O=18-3E-3$.  The first solution gives
$E=6$ and $O=-3$, impossible.  The second gives $E=4$ and $O=3$.
The fibers $B,C$ are entirely even, and $P,Q$ have multiplicity one.
Consequently
\[
 2o_A+o_\Sigma=3,\qquad 0\leq o_A\leq1,\quad
 0\leq o_\Sigma\leq2.
\]
Thus $o_A=o_\Sigma=1$, and
$t_\Sigma=8-2e_\Sigma-2o_\Sigma=2$, contradicting the lemma.

The portable verifier \texttt{verify.py} replays the frozen frontier,
checks the signed identities, and supplies directly checked orbit-count
witnesses for the other $1669$ profiles.  It also independently checks
this two-case proof for profile $133$.  The broader replay imposes
Galois equality of all three $h$-orbit counts, spin-fiber and spin-fixed
four-divisibility, and the prior dimension, rank, involution and
AMV/sixteenth-indicator conditions.  Its receipt is
\texttt{certificate.json}; no optimizer or full center fusion
reconstruction is used.  Passing this necessary test does not assert
existence of spins, modular data, or a categorification.

